\documentclass[10pt,reqno]{amsart}

\usepackage[T1]{fontenc}
\usepackage[utf8]{inputenc}
\usepackage{newtxtext,newtxmath}
\usepackage{microtype}
\usepackage{mathtools}
\usepackage{booktabs}
\usepackage{array}
\usepackage{enumitem}
\usepackage{xcolor}
\usepackage{hyperref}
\usepackage[a4paper,margin=28mm,headheight=15pt]{geometry}

\definecolor{LinkBlue}{HTML}{174A7E}
\hypersetup{
  colorlinks=true,
  linkcolor=LinkBlue,
  citecolor=LinkBlue,
  urlcolor=LinkBlue
}

\setlength{\parskip}{2pt}
\setlist[itemize]{leftmargin=6mm,itemsep=1.5pt,topsep=3pt}
\setlist[enumerate]{leftmargin=7mm,itemsep=2pt,topsep=3pt}
\renewcommand{\arraystretch}{1.15}
\emergencystretch=2em

\newtheorem{theorem}{Theorem}[section]
\newtheorem{proposition}[theorem]{Proposition}
\newtheorem{lemma}[theorem]{Lemma}
\newtheorem{corollary}[theorem]{Corollary}
\theoremstyle{definition}
\newtheorem{definition}[theorem]{Definition}
\newtheorem{remark}[theorem]{Remark}

\title[Counterexamples to Nahm's Conjecture 4.1 on $q$-Hypergeometric Modularity]
{Counterexamples to Nahm's Conjecture on the Bloch-Group Torsion Criterion for Modularity}
\author{Argus Open Research Project}
\date{August 27, 2026}
\subjclass[2020]{Primary 11F03; Secondary 19F27, 33D15}
\begin{document}
\begin{abstract}
We record the counterexamples of Vlasenko and Zwegers (2011) showing that the $q$-hypergeometric Nahm series $f_{A,B,C}(\tau)$ can be modular for positive-definite rational $r\times r$ matrices $A$ even when the associated element in the Bloch group fails the torsion condition of Nahm's Conjecture~4.1.
\end{abstract}
\maketitle
\section{Main Refutation}
\begin{theorem}[Vlasenko--Zwegers, 2011]
For rank $r\ge 2$, there exist explicit positive-definite symmetric matrices $A\in\mathrm{Mat}_r(\mathbb{Q})$, vectors $B\in\mathbb{Q}^r$, and scalars $C\in\mathbb{Q}$ such that the Nahm series
\[
  f_{A,B,C}(\tau) \;=\; \sum_{n\in\mathbb{Z}_{\ge 0}^r} \frac{q^{\frac{1}{2}n^{\mathsf T}An + B^{\mathsf T}n + C}}{(q)_n}
\]
is a modular function, yet the real solution of the Nahm system $1-z_i=\prod_{j=1}^r z_j^{A_{ij}}$ does not define a torsion element in the Bloch group $\mathcal{B}(\overline{\mathbb{Q}})$.
\end{theorem}
\end{document}
